% --- MAGIC COMMENTS: CONFIGURAZIONE EDITOR --- (vedi main.tex per riferimenti e spiegazioni) ---
% !TEX root = ../../main.tex

% _____________________________________________________ BIBLIOGRAFIA ___________________________________________________

\begin{comment}     NEL FILE PACKAGES SONO STATE RIPORTATE LE SEGUENTI ISTRUZIONI, RELATIVE ALLA BIBLIOGRFIA.
% --- PACCHETTI PER BIBLIOGRAFIA E PER GESTIONE SMART DEGLI APICI ---
\usepackage{biblatex}               % Pacchetto per la gestione avanzata della bibliografia, miglior standard attuale
\usepackage{csquotes}               % Richiesto da biblatex per la gestione intelligente delle virgolette e degli apici
                                    % nelle citazioni bibliografiche. Garantisce che le citazioni siano formattate
                                    % correttamente in base alla lingua del documento (es. virgolette italiane « »).

\addbibresource{chapters/00.2_Bibliography/HKN_ElnBiliography.bib}
% File (Database bibliografico) contentnente tutte le fonti bibliografiche di interesse per il progetto.
% ATTENZIONE: il file viene importato nel main.tex con il comando \input. Tieni a mente che TeX è un linguaggio
% Espansore di macro: quindi il codice è come se fosse scritto nel main.tex, e per questo motivo è importante che i
% percorsi siano riferti alla posizione del main.tex.
\end{comment}

% \backmatter inserito prima della chiamata di questo file. Inoltre vedi spiegazioni a riguardo nel commento che segue.

\nocite{*}              % Opzione di biblatex per far stampare in bibliografia tutte le fonti presenti nel database
                        % bibliografico (file .bib)
\printbibliography      % Stampa la bibliografia
